\section{Detecting hallucinations in large language models using semantic entropy}
\paperhead{Detecting hallucinations in large language models using semantic entropy}{Sebastian Farquhar, Jannik Kossen, Lorenz Kuhn, Yarin Gal}{Nature 630(8017), pp. 625 to 630, 2024}{https://doi.org/10.1038/s41586-024-07421-0}{OpenAlex title match 1.00, cited\_by=725}

\subsection{What the paper does}

The paper attacks a narrow but tractable slice of the hallucination problem. It separates \emph{confabulations}, which the authors define as arbitrary and incorrect generations that are sensitive to the random seed, from other failure modes such as errors learned systematically from training data. A confabulation is a generation the model would not repeat if resampled, so it is detectable without any labelled ground truth for the question being asked. The practical requirement is that the detector must work on new and unseen questions to which a human supervisor may not know the answer, which rules out supervised probes trained on task-specific correctness labels.

The core idea is that uncertainty should be measured over \emph{meanings}, not token sequences. Standard predictive entropy over sampled sequences is confounded because many distinct token strings express the same answer, so a model that is completely certain of a fact but flexible about phrasing looks uncertain. Semantic entropy first collapses samples into meaning-equivalence classes and only then computes entropy over those classes, so paraphrase variation cancels out and only genuine disagreement about the answer registers.

\subsection{How the method works}

For a prompt $x$, the method samples $M$ candidate outputs (10 by default) at temperature 1 with nucleus sampling ($P=0.9$) and top-$K$ sampling ($K=50$), plus one low-temperature generation at $0.1$ taken as the best answer. The samples are partitioned into clusters $C_i$ by \emph{bidirectional entailment}: two outputs join the same cluster only if, in the context of the question, each entails the other. Sentence-length experiments use GPT-3.5 as the entailment judge, chosen because its entailment predictions agreed well with human raters; the paragraph-length experiments use DeBERTa-Large-MNLI. Cluster probabilities are aggregated from the token likelihoods of their members, with length normalisation by arithmetic mean log-probability rather than sum so that longer sequences do not inflate the estimate, and the semantic entropy is the entropy over the resulting cluster distribution. A \emph{discrete} variant drops the need for token probabilities entirely by estimating $P(C_i \mid x)$ from cluster frequencies, treating each generated output as equally probable, which makes the method applicable to black-box APIs.

\begin{center}
\begin{tikzpicture}[node distance=1.0cm, >=Stealth, every node/.style={font=\scriptsize}]
\node[draw, rounded corners, text width=2.2cm, align=center, minimum height=9mm] (q) {prompt $x$};
\node[draw, rounded corners, text width=2.4cm, align=center, minimum height=9mm, right=of q] (s) {sample $M{=}10$ outputs, $T{=}1$};
\node[draw, rounded corners, text width=2.6cm, align=center, minimum height=9mm, right=of s] (c) {bidirectional entailment clusters $C_1 \dots C_{|C|}$};
\node[draw, rounded corners, text width=2.4cm, align=center, minimum height=9mm, right=of c] (e) {entropy over $P(C_i \mid x)$};
\node[draw, rounded corners, text width=2.0cm, align=center, minimum height=9mm, right=of e] (g) {reject or answer};
\draw[->] (q) to (s); \draw[->] (s) to (c); \draw[->] (c) to (e); \draw[->] (e) to (g);
\end{tikzpicture}
\end{center}

Paragraph-length generation is handled by decomposition: an LLM splits the paragraph into factual claims, generates $Q$ questions that could have produced each claim, samples $M$ answers per question, computes semantic entropy per question including the original claim, and averages over questions. Baselines are naive sequence entropy, P(True) with few-shot prompting, and a supervised embedding regression on final hidden states.

\subsection{What it reports}

Evaluation covers TriviaQA, SQuAD 1.1, BioASQ, NQ-Open and SVAMP for sentence-length answers, with LLaMA 2 Chat (7B, 13B, 70B), Falcon Instruct (7B, 40B) and Mistral Instruct (7B); FactualBio with GPT-4 covers paragraph-length biography generation. Correctness is graded by GPT-4 against a reference answer, with SQuAD F1 used for short phrases. Averaged over the 30 task and model combinations, semantic entropy reaches the best AUROC of $0.790$, ahead of naive entropy ($0.691$), P(True) ($0.698$) and the embedding regression baseline ($0.687$). The paper also reports AURAC, the area under the rejection-accuracy curve, which measures accuracy on the questions that survive after the most uncertain ones are refused; on the paragraph-length biography task discrete semantic entropy beats a self-check baseline and a P(True) variant. The gains hold across model families and scales, and require no domain knowledge or task-specific labels.

\subsection{How Mytheus uses this}
\begin{takeaway}
Use this as the offline reference standard for the $0.7$ and $0.4$ confidence gates, not as an online detector. Build a golden-set harness stage in CI that computes semantic entropy over $M$ resamples of every \texttt{DecisionAgent} and \texttt{ActionAgent} extraction and reports AUROC of self-reported confidence versus semantic entropy against the role labels, alongside the existing Brier score. If a self-reported number materially underperforms semantic entropy at separating correct from confabulated extractions, the $0.7$ gate is miscalibrated and the paper gives a Nature-published baseline to say so with. The discrete variant is the one to implement: it needs no token log-probabilities from LiteLLM and runs its entailment pass on a locally hosted NLI cross-encoder, which sits naturally beside the BGE reranker on the Ray Serve GPU tier and keeps tenant text off a third-party judge, unlike the paper's GPT-3.5 entailment and GPT-4 grading which would violate the PII-redaction-before-egress rule. Be blunt about where it does not fit: $M{=}10$ samples plus a pairwise entailment pass cannot run inside the \texttt{ContextGapAgent} p95 1.5 second peer round trip, and it multiplies token spend against per-tenant hard quotas, so it must not sit on the live path. The defensible compromise is to run it only on the Opus-routed high-stakes slice (conflict adjudication and end-of-meeting synthesis, under 1 percent of call volume) where the budget is post-meeting, and to reuse the FactualBio decomposition shape at the RAG \texttt{cite} stage: decompose the assembled answer into claims and drop a claim whose per-claim entropy is high even when it binds to a resolvable citation, which catches the case citation binding alone misses. Measuring and reporting exactly where semantic entropy fits and where it blows the latency envelope is itself a publishable contribution.
\end{takeaway}

\subsection{Limitations and what it leaves open}
\begin{itemize}[leftmargin=*, itemsep=1pt, topsep=2pt]
  \item Scoped to confabulations only: the method explicitly does not address models that are confidently wrong because of systematic errors learned in training, so an extraction that is consistently wrong across all resamples (a decision attributed to the wrong stakeholder every time) draws low entropy and passes any gate built on it.
  \item No latency or throughput budget is measured anywhere in the main results: $M{=}10$ generations plus an entailment pass over sample pairs is evaluated purely on AUROC and AURAC, and the cost discussion is confined to supplementary material, leaving the method untested against any hard deadline such as a live agent round trip.
  \item Never tested in a retrieval-augmented setting: uncertainty is computed over free generation from parametric knowledge, so it is unknown whether semantic entropy still separates grounded from ungrounded output once the model is handed retrieved passages under a citation contract, which is exactly the configuration of the Mytheus \texttt{assemble} and \texttt{cite} stages.
  \item Assumes a single shared ground truth per question, graded by GPT-4 against one reference answer or SQuAD F1, so the method can flag that an answer is uncertain but has no notion of whether an answer is correct \emph{for a particular reader}, which is the per-participant relevance judgement \texttt{PersonalListenerAgent} is built on.
  \item Evaluated only on static single-turn benchmarks with clean written prompts: there is no diarized or streaming speech, no mutable interim transcript, no multi-speaker attribution and no multi-tenant setting, so nothing is known about behaviour when the input itself is still being revised or when isolation constrains which samples may be drawn.
\end{itemize}
